\begin{lstlisting}[
  caption={Prompt template for inferring abstraction concepts for next step of retrieval.},
  label={lst:apdx_prompt_get_plan},
]
============================================================
Prompt Get_Plan
============================================================

You are an expert at analyzing an agent's goal and current observation and generating retrieval tags for a goal-directed question-answering system.
Your setting:

* Goal: The agent's overall objective to accomplish.
* Current Subgoal: The subgoal the agent is currently pursuing. (can be None)
* Current State: A description of the agent's current internal state. (can be None)
* Input (Current Observation): The agent's most recent observation (typically a question or task instruction).
* Task: Extract a prioritized set of high-quality tags that are most likely to retrieve information that directly helps accomplish the Goal.

Instructions:

1. Goal-directed Tag Selection (CRITICAL):

* Read the Goal and the Current Observation carefully.
* Only generate tags that are HIGHLY LIKELY to retrieve evidence needed to solve/complete the Goal.
* Prefer tags that identify:
* The target entity/entities the Goal is asking about (people, organizations, places, works, events).
* Bridge entities implied by the observation that are likely required for multi-hop retrieval.
* Explicit constraints: dates, years, roles, titles, unique descriptors, numbers.


* Avoid low-signal or generic tags that are unlikely to retrieve helpful evidence (e.g., "known for", "famous", "character" unless the Goal specifically depends on them).

2. Concrete, Grounded Tags:

* The MAJORITY of tags MUST be short text spans copied VERBATIM from the Goal or the Current Observation (exact substrings).
* You MAY add a SMALL number of non-literal tags only if they are short, strongly implied, and clearly necessary for retrieval (e.g., a canonical name expansion or a standard alias).
* If a tag is a verb, you MUST use its base (lemma) form (e.g., "direct", not "directed" or "directing").

3. Prioritization & Quantity:

* the number of tags in total should be proper. (not "as many as possible").
* Sort tags by expected retrieval usefulness (most useful first).
* Ensure tags are content-bearing and relatively short.

4. CRITICAL -- Forbidden tags:

* Do NOT generate the tag "user".
* Do NOT use meta-labels or type names such as "Name", "Person", "Year", "Date", "City", "Country", "Location", "Genre", "FilmTitle", etc.

Output Format:

### Reasoning

[You process of analyzing the information and completing the task]

### Tags

**Tags:** ["tag0", "tag1", "tag2", "tag3", ...]
(for example: "Central Area", "focal point", "famous for", etc)

### Next Subgoal

## [A single best next subgoal that the agent should pursue now.]

Input:
Goal: {goal}
Current Subgoal: {subgoal}
Current State: {state}
Current Observation: {observation}
\end{lstlisting}
